% ============================================================
% HYPOTHESIS 2 — THEORY
% ============================================================

\section{Hypothesis 2 — Verification}

\begin{frame}{Hypothesis 2: Theory}

\textbf{Hypothesis 2:}

Nếu LM loss thay đổi theo retrieval coefficient $d$, gradient

\[
\frac{\partial L_{\mathrm{LM}}}{\partial d}
\]

có thể cung cấp tín hiệu cục bộ cho retrieval decision.

\vspace{0.25cm}

\textbf{Gradient interpretation:}

\begin{itemize}
    \item
    \[
    \frac{\partial L_{\mathrm{LM}}}{\partial d}<0
    \]
    Tăng $d$ có xu hướng làm giảm LM loss.

    \item
    \[
    \frac{\partial L_{\mathrm{LM}}}{\partial d}>0
    \]
    Tăng $d$ có xu hướng làm tăng LM loss.
\end{itemize}

\vspace{0.2cm}

\textbf{Verification:}

Với mỗi decision level:

\[
d=0/32,\ldots,32/32
\]

tính:

\[
\frac{\partial L_{\mathrm{LM}}}{\partial d}
\]

\end{frame}


% ============================================================
% HYPOTHESIS 2 — EXAMPLE
% ============================================================

\begin{frame}{Hypothesis 2: Example}

\textbf{Question:}

\textit{Who wrote Common Sense and why was it important?}

\vspace{0.2cm}

\textbf{Retrieved Documents:}

\begin{itemize}
    \item \textbf{Common Sense:}

    Common Sense was a pamphlet written by Thomas Paine.

    \item \textbf{Thomas Paine:}

    Thomas Paine was an author and key figure in the American Revolution. His work Common Sense crystallized the rebellious demand for independence from Great Britain.
\end{itemize}

\vspace{0.15cm}

\textbf{Experiment:}

Giữ cố định question, answer và documents.

\vspace{0.1cm}

Thay đổi:

\[
d=0/32,\ldots,32/32
\]

\vspace{0.1cm}

Với mỗi $d$, tính:

\[
L_{\mathrm{LM}}(d)
\]

và:

\[
\frac{\partial L_{\mathrm{LM}}}{\partial d}.
\]

\vspace{0.15cm}

\textbf{Expected observation:}

\begin{itemize}
    \item $\frac{\partial L_{\mathrm{LM}}}{\partial d}<0$: retrieved information có ích.
    \item $\frac{\partial L_{\mathrm{LM}}}{\partial d}>0$: retrieved information không cần thiết hoặc gây hại.
\end{itemize}

\end{frame}